% !TEX root = ../main.tex
\section{Introduction}
\label{sec:introduction}

High-speed printed circuit board (PCB) interconnects range from planar traces to complex
multilayer routing structures. Their electrical behaviour varies
with frequency and can significantly degrade signal integrity. These effects are
described by scattering parameters (S-parameters) to evaluate signal-integrity (SI),
which quantify insertion and return losses as well as crosstalk
\cite{pozar2011microwave}. Electromagnetic simulation can accurately capture this
behaviour. But when there are many geometric and material variations, the evaluation
becomes computationally expensive \cite{bruns2007numerical}. Machine learning offer a
faster alternative by learning the relationship between design parameters and simulated
responses \cite{swaminathan2020demystifying,akinwande2024complex}.

This work follows a staged workflow. Linear and non-linear baselines establish reference
performance. Then, neural models learn broadband response as a continuous curve. A
frequency-conditioned network predicts the reciprocal complex 12-port S-matrix from ten
design parameters across 200 frequency points. Evaluation on designs held out from
training considers predictive error, reciprocity, passivity and difficult response
regions. The selected full-matrix model improves upon a training-mean complex reference
and guarantees reciprocity. It predicts a broader response than the selected six-curve
model but has higher $S_{7,1}$ insertion-loss (IL) error. Since their architectures and
targets also differ, the comparison does not attribute this difference solely to output
scope.
